\subsection{曲面情形}\label{sec:20c}
剩下的 $n=1$ 情形更为棘手。我们将采用一种较为间接的方法，把显式计算与抽象分类理论结合起来。以下始终假设系数域满足 $\operatorname{char}(\K)=0$；这一假设会被反复用到，甚至在论证中纯代数的部分也是如此。例如，若取 $\operatorname{char}(\K)=2$，\stref{引理}{20.2}就不成立。

首先，我们考察一种可能导致上同调为 $A=A_m^1$ 的 $\Ai$ 范畴 $\A$ 不具形式性的情形。不失一般性，可设 $\mu_{\A}^1=0$，这样 $\A$ 就是 $A$ 本身上的一个 $\Ai$ 结构。还可假设 $\A$ 严格含幺。
\SourcePage{316}
\StatementNumber{1}\begin{lemma}\label{st:20.1}
再假设 $\mu_{\A}^3=0$，且 $m\ge5$。考虑
\begin{equation}\label{eq:20.8}
\begin{aligned}
&\mu_{\A}^4\bigl((3,2),(2,1),(1,2),(2,3)\bigr)
 -\mu_{\A}^4\bigl((3,2),(2,3),(3,2),(2,3)\bigr)\\
&\quad+\mu_{\A}^4\bigl((3,4),(4,3),(3,2),(2,3)\bigr)
 +\mu_{\A}^4\bigl((3,2),(2,3),(3,4),(4,3)\bigr)\\
&\quad-\mu_{\A}^4\bigl((3,4),(4,3),(3,4),(4,3)\bigr)
 +\mu_{\A}^4\bigl((3,4),(4,5),(5,4),(4,3)\bigr)=q(3),
\end{aligned}
\end{equation}
其中 $q\in\K$。若 $q\ne0$，则 $\A$ 不是形式的。
\end{lemma}
\begin{proof}
设 $\A,\widetilde{\A}$ 为两个这样的 $\Ai$ 结构，并以 $q,\widetilde q$ 表示它们按\eref{20.8}得到的系数。假设有一个严格含幺的形式微分同胚 $\Phi:\A\to\widetilde{\A}$，其一次项为 $\Phi^1=\Id$。直接计算可得
\begin{equation}\label{eq:20.9}
(\widetilde q-q)(3)=
\bigl(\Phi^2\bigl((3,2),(2,3)\bigr)-\Phi^2\bigl((3,4),(4,3)\bigr)\bigr)^2.
\end{equation}
再考虑其逆形式微分同胚 $\Psi$，它满足 $\Psi^2=-\Phi^2$，并考察 $q-\widetilde q$ 对应的等式。比较两式可见，\eref{20.9}的右端为零，因而 $q=\widetilde q$。\znote{308}{两处“显然”值得展开。其一，$m\ge5$ 是必需的：(20.8) 中出现 $(3,4),(4,5),(5,4),(4,3)$，用到顶点 $5$，故 $m\ge5$ 时六项才都有意义。其二，(20.9) 的右端是平方，因此把同一式用于逆微分同胚便得到 $x^2=-x^2$，特征不为 $2$ 时 $x=0$；即 $q$ 在形式同构下不变，$q\ne0$ 就意味着非形式。注意 (20.9) 只用到 $\mu^3=0$ 与 $\Phi^1=\Id$ 两个条件，故它对任意 $m\ge5$ 的 $\A$ 都成立；“$q\ne0$ 则不存在到平凡结构的形式微分同胚”正是非形式的定义。[来源：A，原书引理 20.1 及其证明。]}

现在将此结论应用于给定的 $\A$，并假设它是形式的。那么，可以找到上述类型的形式微分同胚 $\Phi$，使 $\widetilde{\A}$ 为平凡的 $\Ai$ 结构，因此 $q$ 必须为零。
\end{proof}
